% Generated draft: requires mathematical and source review.
\documentclass[11pt]{article}
\usepackage[margin=25mm]{geometry}
\usepackage{fontspec}
\setmainfont{Latin Modern Roman}
\usepackage{amsmath,amssymb,mathtools,unicode-math}
\setmathfont{Latin Modern Math}
\usepackage{enumitem,needspace,xurl,hyperref}
\setlength{\parindent}{0pt}
\setlength{\parskip}{5pt}
\setlength{\emergencystretch}{3em}
\newcommand{\N}{\mathbb{N}}
\newcommand{\Z}{\mathbb{Z}}
\newcommand{\Q}{\mathbb{Q}}
\newcommand{\R}{\mathbb{R}}
\newcommand{\C}{\mathbb{C}}
\newcommand{\F}{\mathbb{F}}
\newcommand{\A}{\mathbb{A}}
\newcommand{\PP}{\mathbb{P}}
\newcommand{\E}{\mathbb{E}}
\newcommand{\eps}{\varepsilon}
\newcommand{\dd}{\,\mathrm{d}}
\begin{document}
\section*{UnsolvedMath ID: 6700028}
\textbf{Source problem code:} catalogue entry
\subsection{Definitions and mathematical statement}
The source says that \(X\) is ``homomorphic to \(\mathbb{CP}^2\)'', which is not a specified topological equivalence of manifolds. The likely homeomorphism reading must be checked; it is not strengthened to diffeomorphism. In the natural metric setting, let \((X,g)\) be a smooth Riemannian four-manifold with the intended topological relation to \(\mathbb{CP}^2\), and for a surface \(S\subset X\) put
\[
 U_\delta(S)=\{x\in X:d_g(x,S)<\delta\}.
\]
The requested numerical implication is
\[
 \left[\ \forall\ T(X)\text{-non-spin surfaces }S\subset X\ 
 \forall\delta\in(0,1]:\quad
 \operatorname{Vol}_g(U_\delta(S))\ge\delta^2\ \right]
 \quad\Longrightarrow\quad
 \operatorname{spin\text{-}area}(X)\ge10^{-6}.
\]
Both the universal surface condition and every real radius \(0<\delta\le1\) are retained; the constant is exactly \(1/1{,}000{,}000\).

The excerpt does not define \(T(X)\)-non-spin. One plausible tangent-bundle reading is that the pullback \(i^*TX\), \(i:S\hookrightarrow X\), has no spin structure; for an oriented ambient tangent bundle this is detected by \(i^*w_2(TX)\ne0\). That differs from saying the surface's own tangent bundle is non-spin. The source does not specify the surface regularity, closedness, embedding versus mapping, or whether this is the intended obstruction criterion.

Most importantly, \(\operatorname{spin\text{-}area}\) and its normalization are not defined in the supplied record. It is not equated with a particular K-area invariant merely because the summary mentions that program. The metric convention and the meaning of the fixed scale \(1\) are also essential for the numerical assertion. The source summary calls the estimate open and says no published verification of this specific \(\mathbb{CP}^2\)/\(10^{-6}\) claim was found; that is a supplied review report, not independent current verification.
\subsection{Short English statement}
If every surface detecting the intended non-spin tangent-bundle obstruction has tube volume at least \(\delta^2\) for all \(0<\delta\le1\), must the spin-area be at least \(10^{-6}\)? The topological wording and both specialized invariant definitions need confirmation from the original source.
\subsection{Sources}
\begin{itemize}
\item Local registry: \url{data/unsolvedmath/problems.json}, permanent UnsolvedMath ID \(6700028\), source code catalogue entry; supplied statement, background, and research summary, updated 2026-08-06.
\item Misha Gromov, \emph{101 Questions, Problems and Conjectures around Scalar Curvature} (2017), printed tag [?28], occurrence 29, PDF page 38: \url{https://www.ihes.fr/~gromov/wp-content/uploads/2018/08/101-problemsOct1-2017.pdf}. Supplied normalized PDF-text extraction accessed 2026-07-29; described as incomplete.
\item No separate definition of spin-area, tangent-bundle-non-spin surfaces, or publication verifying this quantitative estimate is supplied.
\end{itemize}
\end{document}
